\documentclass[10pt,a4paper]{amsart}
\usepackage[margin=19mm]{geometry}
\usepackage{amsmath,amssymb,mathtools}
\usepackage[scr=rsfs,cal=euler]{mathalfa}
\usepackage{xfrac}
\usepackage[unicode,hidelinks,bookmarks=false]{hyperref}
\newcommand{\ie}{i.e.\ }
\newcommand{\cf}{cf.\ }
\newcommand{\resp}{respectively\ }
\newcommand{\Subsection}{\subsection}
\providecommand{\nobreakdash}{\nobreak\hbox{-}}
\setlength{\emergencystretch}{2em}
\multlinegap0pt
\usepackage{booktabs}
\newtheorem{theorem}{Theorem}
\newtheorem{lemma}[theorem]{Lemma}
\newtheorem{proposition}[theorem]{Proposition}
\newtheorem{corollary}[theorem]{Corollary}
\newtheorem{conjecture}[theorem]{Conjecture}
\newtheorem{remark}[theorem]{Remark}
\newcommand{\F}{\mathbb{F}}
\newcommand{\HHull}{\mathrm{Hull}_H}
\newcommand{\Aut}{\mathrm{Aut}}
\begin{document}
\title[Comparative monotonicity of linear codes by Hermitian and sy]{Comparative monotonicity of linear codes by Hermitian and symplectic hull dimensions}
\author{Keita Ishizuka}
\date{}
\maketitle
\noindent\emph{Source and licence.} Comparative monotonicity of linear codes by Hermitian and symplectic hull dimensions. Version arXiv:2605.29204v1 (CC BY 4.0). This material is distributed under the Creative Commons Attribution 4.0 International licence, http://creativecommons.org/licenses/by/4.0/, as recorded in the arXiv OAI licence metadata for this exact version and shown on the abstract page https://arxiv.org/abs/2605.29204v1. Full rights notice: licenses/arxiv-2605.29204-CC-BY-4.0.txt.\par\smallskip
\noindent\textit{Editorial scope.} Counted unit: the original \texttt{theorem} environment \texttt{thm:main} at source lines 411--428 together with its complete proof at lines 440--518; the delivered document keeps source lines 123--948 so that every referenced label and every citation resolves inside it. Native editable source is the paper's own concatenated TeX; the AMS wrapper is editorial formatting only. No publisher class, upstream script or unrelated artwork is redistributed.
\section{Preliminaries} \label{sec:prelim}
%% ============================================================

\subsection{Linear codes and generator matrices}\label{sec:prelim-codes}

Let $\F_q$ denote the finite field with $q$ elements, where $q$ is a prime power.
A \emph{linear $[n, k]_q$ code} $C$ is a $k$-dimensional subspace of $\F_q^n$.
The integer $n$ is the \emph{length}, $k$ is the \emph{dimension}, and the elements of $C$
are called \emph{codewords}. A \emph{generator matrix} of $C$ is a $k \times n$ matrix
$G$ over $\F_q$ whose rows form a basis of $C$; equivalently, $C = \{\mathbf{u} G : \mathbf{u} \in \F_q^k\}$.
The \emph{Hamming weight} of a codeword $\mathbf{c} \in C$ is the number of nonzero coordinates,
and the \emph{minimum distance} of $C$ is $d(C) = \min\{\mathrm{wt}(\mathbf{c}) : \mathbf{c} \in C \setminus \{\mathbf{0}\}\}$.
We sometimes write $[n, k, d]_q$ for an $[n, k]_q$ code of minimum distance $d$.

\subsection{Bilinear forms and dual codes}\label{sec:prelim-dual}

A \emph{bilinear form} on $\F_q^n$ is a $\F_q$-bilinear map $B : \F_q^n \times \F_q^n \to \F_q$.
Three standard examples will concern us:
\begin{itemize}
\item the \emph{Euclidean inner product} $\langle \mathbf{x}, \mathbf{y} \rangle = \sum_{i=1}^n x_i y_i$
      on $\F_q^n$, which is symmetric;
\item the \emph{Hermitian inner product} $\langle \mathbf{x}, \mathbf{y} \rangle_H = \sum_{i=1}^n x_i y_i^q$
      on $\F_{q^2}^n$, which is sesquilinear with respect to the Frobenius automorphism $x \mapsto x^q$;
\item the \emph{symplectic inner product} on $\F_q^{2n}$ defined by
      \[
        \langle \mathbf{x}, \mathbf{y} \rangle_S = \sum_{i=1}^{n} (x_i y_{n+i} - x_{n+i} y_i),
      \]
      which is alternating: $\langle \mathbf{x}, \mathbf{x} \rangle_S = 0$ for all $\mathbf{x}$.
      Equivalently, the symplectic Gram matrix is $\Omega_n = \begin{pmatrix} O & I_n \\ -I_n & O \end{pmatrix}$.
\end{itemize}
For a linear code $C$ over $\F_q$ (resp.\ $\F_{q^2}$), the \emph{Euclidean dual},
\emph{Hermitian dual}, and \emph{symplectic dual} are defined respectively by
\begin{align*}
  C^\perp &= \{\mathbf{y} \in \F_q^n : \langle \mathbf{x}, \mathbf{y} \rangle = 0 \ \forall \mathbf{x} \in C\}, \\
  C^{\perp_H} &= \{\mathbf{y} \in \F_{q^2}^n : \langle \mathbf{x}, \mathbf{y} \rangle_H = 0 \ \forall \mathbf{x} \in C\}, \\
  C^{\perp_S} &= \{\mathbf{y} \in \F_q^{2n} : \langle \mathbf{x}, \mathbf{y} \rangle_S = 0 \ \forall \mathbf{x} \in C\},
\end{align*}
where the symplectic dual is defined for an $[2n, k]_q$ code $C \subseteq \F_q^{2n}$.
Each dual code is itself a linear code of dimension $n - k$ (resp.\ $2n - k$ in the symplectic case),
since all three forms are non-degenerate.
For a generator matrix $G$ of $C$, the dual is generated by any matrix $H$ satisfying
$G H^T = O$ (Euclidean), $G \overline{H}^T = O$ (Hermitian), or $G \Omega_n H^T = O$ (symplectic),
where $\overline{H} = (h_{ij}^q)$ is the entrywise $q$-th power
and $\Omega_n$ is the symplectic Gram matrix.

\subsection{Hull, LCD codes, and self-orthogonal codes}\label{sec:prelim-hull}

For each of the three inner products, the \emph{hull} is defined as the intersection of the code
with its dual:
\begin{align*}
  \mathrm{Hull}(C) &:= C \cap C^\perp & &\text{(Euclidean),}\\
  \HHull(C) &:= C \cap C^{\perp_H} & &\text{(Hermitian),}\\
  \mathrm{Hull}_S(C) &:= C \cap C^{\perp_S} & &\text{(symplectic).}
\end{align*}
The hull dimension interpolates between two extremes: $C$ is called a
\emph{linear complementary dual (LCD)} code if its hull is trivial
(introduced by Massey~\cite{Massey1992}), and \emph{self-orthogonal} if $C$
is contained in the appropriate dual.
Binary LCD codes have applications to side-channel attack countermeasures~\cite{CarletGuilley2016}
and decoding complexity, while self-orthogonal codes underlie the construction of
quantum error-correcting codes via the stabilizer and entanglement-assisted frameworks~\cite{KKKS2006,Dcc-EAqecc}.

The hull dimension can be read off from any generator matrix:
\begin{proposition}\label{prop:hull-dim}
Let $C$ be a linear code with generator matrix $G$.
\begin{enumerate}
\item \cite[Proposition 3.2]{Dcc-EAqecc} If $C$ is an $[n, k]_{q^2}$ code, then
$\dim(\HHull(C)) = k - \mathrm{rank}(G \overline{G}^T)$.
In particular, $C$ is Hermitian LCD if and only if $G \overline{G}^T$ is nonsingular,
and Hermitian self-orthogonal if and only if $G \overline{G}^T = O$.
\item If $C$ is an $[2n, k]_q$ code, then
$\dim(\mathrm{Hull}_S(C)) = k - \mathrm{rank}(G \Omega_n G^T)$.
In particular, $C$ is symplectic LCD if and only if $G \Omega_n G^T$ is nonsingular,
and symplectic self-orthogonal if and only if $G \Omega_n G^T = O$. Since $G \Omega_n G^T$
is alternating, $\mathrm{rank}(G \Omega_n G^T)$ is even, hence
$k - \dim(\mathrm{Hull}_S(C))$ is even.
\end{enumerate}
\end{proposition}
For Hermitian, $0 \leq \dim(\HHull(C)) \leq \min(k, n-k)$.
For symplectic, $0 \leq \dim(\mathrm{Hull}_S(C)) \leq \min(k, 2n-k)$ with the
additional constraint that $k - \dim(\mathrm{Hull}_S(C))$ is even.

\subsection{The Euclidean prototype: Bouyuklieva-Bouyukliev-\"Ozbudak}\label{sec:prelim-bbo}

The starting point of our investigation is the following result of Bouyuklieva,
Bouyukliev, and \"Ozbudak~\cite{BBOz2026}, which gave the first systematic study of
how the count $A_{n,k,\ell,q}$ of $[n,k]_q$ Euclidean codes with hull dimension $\ell$
varies with $\ell$.

\begin{theorem}[\cite{BBOz2026}, Lemma~8]\label{thm:BBO-odd}
Let $q$ be a power of an odd prime, $n$ and $k$ be positive integers with
$k \leq n/2$, and $\ell$ be an integer such that $0 \leq \ell \leq k - 1$. Then
\[
  A_{n,k,\ell,q} = \alpha_{n,k,\ell,q}\,(q^{\ell+1} - 1)\,A_{n,k,\ell+1,q},
\]
where
\[
  \alpha_{n,k,\ell,q} = \begin{cases}
    \displaystyle\frac{q^{n/2-1}}{q^{n/2-1} + \eta((-1)^{n/2})\,q^\ell}
      & \text{if $n$ is even, $k - \ell$ is odd}, \\[8pt]
    \displaystyle\frac{q^{n-k-\ell}}{q^{n-k-\ell} - 1}
      & \text{if $n$ is odd, $k - \ell$ is odd}, \\[8pt]
    \displaystyle\frac{q^{k-\ell}}{q^{k-\ell} - 1}
      & \text{if $n$ is odd, $k - \ell$ is even}, \\[8pt]
    \displaystyle\frac{q^{n/2-\ell}\bigl(q^{n/2-\ell} + \eta((-1)^{n/2})\bigr)}
                       {(q^{n-k-\ell}-1)(q^{k-\ell}-1)}
      & \text{if $n$ is even, $k - \ell$ is even}.
  \end{cases}
\]
If $k - \ell$ is odd, $n \equiv 0 \pmod 4$ or
$n \equiv 2 \pmod 4$ and $q \equiv 1 \pmod 4$, then
$\alpha_{n,k,\ell,q} \geq \tfrac{1}{2}$ with equality when $k = n/2$ and
$\ell = k - 1$. In all other cases $\alpha_{n,k,\ell,q} > 1$.
\end{theorem}

\begin{theorem}[\cite{BBOz2026}, Lemma~10]\label{thm:BBO-even}
Let $q$ be a power of $2$, $n$ and $k$ be positive integers with $k \leq n/2$,
and $\ell$ be an integer such that $0 \leq \ell \leq k - 1$. Then
\[
  A_{n,k,\ell,q} = \alpha_{n,k,\ell,q}\,(q^{\ell+1} - 1)\,A_{n,k,\ell+1,q},
\]
where
\[
  \alpha_{n,k,\ell,q} = \begin{cases}
    \displaystyle\frac{q^{n-\ell-1}}{q^{n-\ell-1} - 1}
      & \text{if $n$ is even, $k - \ell$ is odd}, \\[8pt]
    \displaystyle\frac{q^{n-k-\ell}}{q^{n-k-\ell} - 1}
      & \text{if $n$ is odd, $k - \ell$ is odd}, \\[8pt]
    \displaystyle\frac{q^{k-\ell}}{q^{k-\ell} - 1}
      & \text{if $n$ is odd, $k - \ell$ is even}, \\[8pt]
    \displaystyle\frac{q^{n-\ell}-1}{q^\ell(q^{n-k-\ell}-1)(q^{k-\ell}-1)}
      & \text{if $n$ is even, $k - \ell$ is even}.
  \end{cases}
\]
In all cases $\alpha_{n,k,\ell,q} > 1$.
\end{theorem}

These four-case decompositions reflect the Witt classification of the standard
Euclidean form on $\F_q^n$, in particular the quadratic-residue character
$\eta((-1)^{n/2})$ for $n$ even and $q$ odd. The case $q \equiv 3 \pmod 4$,
$n \equiv 2 \pmod 4$, $k - \ell$ odd is the unique regime where the natural
inequality $A_{n,k,\ell,q} > (q^{\ell+1}-1) A_{n,k,\ell+1,q}$ may fail; in
that case BBO show $A_{n,k,\ell,q} \geq \tfrac{1}{2}(q^{\ell+1}-1) A_{n,k,\ell+1,q}$
(\cite[Theorem~12]{BBOz2026}).

In this paper, we aim to extend
Theorems~\ref{thm:BBO-odd} and~\ref{thm:BBO-even} to the Hermitian
(Section~\ref{sec:main}) and symplectic (Section~\ref{sec:symp}) settings,
in each case deriving a closed-form expression for the ratio factor and an
exact characterization of the parameter set where the BBO-style inequality fails.
In Section~\ref{sec:comp-app} we juxtapose the three results, observing
that the qualitative behavior of hull-monotonicity differs markedly across
the three classical inner products and tracing this difference to the orbit
structure of the corresponding classical group action.

\subsection{Counting linear codes by hull dimension}\label{sec:prelim-counts}

For the Euclidean, Hermitian, and symplectic inner products, we write
\begin{align*}
  A_{n,k,\ell,q}    &= |\{C \subseteq \F_q^n : \dim C = k,\ \dim \mathrm{Hull}(C) = \ell\}|, \\
  A^H_{n,k,\ell,q}  &= |\{C \subseteq \F_{q^2}^n : \dim C = k,\ \dim \HHull(C) = \ell\}|, \\
  A^S_{2n,k,\ell,q} &= |\{C \subseteq \F_q^{2n} : \dim C = k,\ \dim \mathrm{Hull}_S(C) = \ell\}|.
\end{align*}
These are the central objects of our study.

\subsection{Mass formulas for Hermitian and symplectic hull dimension}\label{sec:prelim-mass}

For a prime power $Q$ and integers $0 \leq k \leq n$, the
\emph{Gaussian binomial coefficient}
\[
  \begin{bmatrix} n \\ k \end{bmatrix}_Q
  := \prod_{i=0}^{k-1} \frac{Q^{n-i} - 1}{Q^{i+1} - 1}
\]
counts the number of $k$-dimensional subspaces of an $n$-dimensional vector
space over $\F_Q$. It is a polynomial in $Q$ with non-negative integer
coefficients and satisfies $\binom{n}{k}_Q = \binom{n}{n-k}_Q$. We use this
quantity below as the total count of $[n,k]_Q$ codes against which mass
formulas are verified.

The exact value of $A^H_{n,k,\ell,q}$ was determined by Li et al.\
via classical group action on the set of LCD codes.

\begin{theorem}[{\cite[Theorem~3.11]{LSL2024}}]\label{thm:LSL}
Let $n,k,\ell$ be positive integers with $\ell \leq k \leq n - \ell$.
Let $k_0 = k - \ell$.
Then
\[
  A^H_{n,k,\ell,q} = \begin{cases}
    \displaystyle\left(\prod_{i=1}^{\ell} A_{n,k_0,i}\right) \cdot L(n,k_0,q), & \text{if } n - k_0 \text{ is odd}, \\[10pt]
    \displaystyle\left(\prod_{i=1}^{\ell} B_{n,k_0,i}\right) \cdot L(n,k_0,q), & \text{if } n - k_0 \text{ is even},
  \end{cases}
\]
where
\[
  L(n,k_0,q) = q^{k_0(n-k_0)} \prod_{j=1}^{k_0} \frac{q^{n-k_0+j} - (-1)^{n-k_0+j}}{q^j - (-1)^j}
\]
is the number of Hermitian LCD $[n,k_0]_{q^2}$ codes, and
\begin{align*}
  A_{n,k_0,i} &= \frac{(q^{n-k_0-2i+2}+1)(q^{n-k_0-2i+1}-1)}{q^{2k_0}(q^{2i}-1)}, \\
  B_{n,k_0,i} &= \frac{(q^{n-k_0-2i+2}-1)(q^{n-k_0-2i+1}+1)}{q^{2k_0}(q^{2i}-1)}.
\end{align*}
\end{theorem}

The same authors also obtained a closed mass formula for the symplectic case,
again via classical group action on symplectic LCD codes.

\begin{theorem}[{\cite[Theorem 4.9]{LSL2024}}]\label{thm:LSL-symp}
Let $n, k, \ell$ be positive integers with $\ell \leq k \leq 2n$ and
$k - \ell$ even, written as $k - \ell = 2k_0$. Then
\[
  A^S_{2n,k,\ell,q} = q^{2k_0(n-k_0-\ell)} \cdot
  \frac{\prod_{m=1}^{\ell}(q^{2(n-k_0-\ell+m)}-1)}{\prod_{m=1}^{\ell}(q^m-1)} \cdot
  \begin{bmatrix} n \\ k_0 \end{bmatrix}_{q^2}.
\]
In particular, the number of symplectic LCD $[2n, 2k_0]_q$ codes is
\[
  L_S(n, k_0, q) := A^S_{2n, 2k_0, 0, q} = q^{2k_0(n-k_0)} \begin{bmatrix} n \\ k_0 \end{bmatrix}_{q^2},
\]
and $A^S_{2n, k, \ell, q} = 0$ unless $k - \ell$ is even.
\end{theorem}

\begin{remark}
The printed statement of~\cite[Theorem 4.9]{LSL2024} contains typographical
errors in the product factor: as written, the numerator and denominator exponents
do not depend on the running index in the manner required by the iterative
construction, and the resulting expression fails to yield integer values
(e.g., $7/3$ for $(n,k,\ell,q) = (2,2,2,2)$ instead of the correct $15$).
The form we state above is algebraically equivalent to the recurrence derived in
the proof of~\cite[Theorem 4.9]{LSL2024}, namely
$\prod_{i=1}^{\ell} (q^{2n-2k_0-2i+2}-1)/(q^{2k_0+i}-q^{2k_0}) \cdot
q^{2k_0(n-k_0)} \binom{n}{k_0}_{q^2}$,
obtained by factoring $q^{2k_0+i}-q^{2k_0} = q^{2k_0}(q^i-1)$ and reindexing.
We have verified numerically that this formula satisfies the consistency check
$\sum_\ell A^S_{2n,k,\ell,q} = \binom{2n}{k}_q$ for all parameter ranges
considered in this paper.
\end{remark}

%% ============================================================
\section{Ratio decomposition for Hermitian hull}\label{sec:main}
%% ============================================================

The main result of this section is a closed-form expression for the ratio
$A^H_{n,k,\ell,q} / A^H_{n,k,\ell+1,q}$ together with a sharp characterization
of the parameter set where the natural BBO-style inequality fails.

\subsection{A unified form of the mass formula}\label{sec:unified}

Theorem~\ref{thm:LSL} expresses $A^H_{n,k,\ell,q}$ as one of two products
according to the parity of $s = n - k_0$, with $A_{n,k_0,i}$ used for $s$ odd
and $B_{n,k_0,i}$ for $s$ even. The two cases differ only in the signs in the
numerator factors. We exploit this near-symmetry by introducing a single
sign parameter that uniformly handles both cases. This unified form is the
key technical device for the proof of Theorem~\ref{thm:main} below; it
reduces the four-case analysis that would arise from separately treating
parity transitions to a single algebraic manipulation.

\begin{lemma}[Unified mass formula]\label{lem:unified}
Let $n, k, \ell$ be positive integers with $\ell \leq k \leq n - \ell$, set
$k_0 = k - \ell$, $s = n - k_0$, and $\varepsilon = (-1)^{s+1}$. Define
\begin{equation}\label{eq:Fi}
  F_i(n, k_0, q) := \frac{(q^{s-2i+2} + \varepsilon)(q^{s-2i+1} - \varepsilon)}{q^{2k_0}(q^{2i}-1)}.
\end{equation}
Then for $s$ odd, $\varepsilon = 1$ and $F_i = A_{n,k_0,i}$; for $s$ even,
$\varepsilon = -1$ and $F_i = B_{n,k_0,i}$. Consequently, in both parity cases,
\begin{equation}\label{eq:unified}
  A^H_{n,k,\ell,q} = \left(\prod_{i=1}^{\ell} F_i(n,k_0,q)\right) \cdot L(n,k_0,q).
\end{equation}
\end{lemma}

\begin{proof}
For $s$ odd, $(-1)^{s+1} = (-1)^{\text{even}} = 1$, hence $\varepsilon = 1$ and
$F_i = (q^{s-2i+2}+1)(q^{s-2i+1}-1)/[q^{2k_0}(q^{2i}-1)] = A_{n,k_0,i}$.
For $s$ even, $\varepsilon = -1$ and analogously $F_i = B_{n,k_0,i}$.
Identity~\eqref{eq:unified} then follows directly from Theorem~\ref{thm:LSL}.
\end{proof}

\begin{remark}
The introduction of $\varepsilon$ in Lemma~\ref{lem:unified} simplifies the
parity bookkeeping: passing from $k_0$ to $k_0 - 1$ (which occurs when
incrementing $\ell$ in the ratio decomposition) sends $s \to s+1$ and
$\varepsilon \to -\varepsilon$, so a single sign-flip captures the parity
transition. This is the essential feature that allows the four-case analysis
implicit in Theorem~\ref{thm:LSL} to be replaced by a uniform telescoping
argument in the proof of Theorem~\ref{thm:main}.
\end{remark}

\subsection{Main theorem}\label{sec:Hmain}


\begin{theorem}\label{thm:main}
Let $n, k, \ell$ be positive integers with $\ell+1 \leq k \leq n-\ell-1$,
and set $a = k - \ell$, $b = n-k-\ell$. Then
\begin{equation}\label{eq:decomp}
  A^H_{n,k,\ell,q} = \alpha^H_{n,k,\ell,q}\cdot (q^{\ell+1}-1)\cdot A^H_{n,k,\ell+1,q},
\end{equation}
where
\begin{equation}\label{eq:alpha}
  \alpha^H_{n,k,\ell,q} = \frac{q^{n-2\ell-1}(q^{\ell+1}+1)}{(q^{a}-(-1)^{a})(q^{b}-(-1)^{b})}.
\end{equation}
Moreover, $\alpha^H_{n,k,\ell,q} > 1$ for all valid parameters \emph{except}
when $\ell = 0$, both $a$ and $b$ are odd, and $\min(a,b) = 1$
(equivalently, $\ell=0$, $n$ is even, and $k\in\{1, n-1\}$); in that exceptional case
\[
  \alpha^H_{n,k,0,q} = \frac{q^{\max(a,b)}}{q^{\max(a,b)}+1} \in \left[\tfrac{q}{q+1},\, 1\right).
\]
In all cases, $\alpha^H_{n,k,\ell,q} \geq q/(q+1) \geq 2/3$.
\end{theorem}


\begin{remark}
The lower bound $\alpha^H \geq q/(q+1) \geq 2/3$ is strictly stronger
than the Euclidean bound $\alpha \geq 1/2$ of~\cite{BBOz2026}.
Furthermore, the Hermitian exceptional set is substantially smaller than its
Euclidean counterpart: the latter involves residues modulo~$4$ and quadratic
residuacity in $\F_q^*$, while the former is confined to the boundary
$k\in\{1,n-1\}$ at $\ell = 0$. In the exceptional case, $\alpha^H \to 1$
as $\max(a,b) \to \infty$, so the deficit vanishes asymptotically.
\end{remark}


\begin{proof}[Proof of Theorem~\ref{thm:main}]
By the unified mass formula~\eqref{eq:unified}, the ratio in~\eqref{eq:decomp}
expands as
\[
  \frac{A^H_{n,k,\ell,q}}{A^H_{n,k,\ell+1,q}}
  = \frac{\prod_{i=1}^{\ell}\frac{F_i(n,k_0,q)}{F_i(n,k_0-1,q)}\cdot
    \frac{L(n,k_0,q)}{L(n,k_0-1,q)}}{F_{\ell+1}(n,k_0-1,q)}.
\]
Passing from $k_0$ to $k_0-1$ sends $s\to s+1$, hence $\varepsilon\to -\varepsilon$.
Substituting into the unified $F_i$ definition~\eqref{eq:Fi} and simplifying,
\[
  \frac{F_i(n,k_0,q)}{F_i(n,k_0-1,q)} = \frac{q^{s-2i+1}-\varepsilon}{q^{2}(q^{s-2i+3}-\varepsilon)},
\]
whose product over $i = 1, \ldots, \ell$ telescopes to
\[
  \prod_{i=1}^{\ell}\frac{F_i(n,k_0,q)}{F_i(n,k_0-1,q)}
  = q^{-2\ell} \cdot \frac{q^{s-2\ell+1}-\varepsilon}{q^{s+1}-\varepsilon}.
\]
Direct manipulation of $L(n, k_0, q)$ from Theorem~\ref{thm:LSL} yields
\[
  \frac{L(n,k_0,q)}{L(n,k_0-1,q)} = q^{s-k_0+1}\cdot\frac{q^{s+1}-\varepsilon}{q^{k_0}-(-1)^{k_0}},
\]
and a straightforward substitution gives
\[
  F_{\ell+1}(n,k_0-1,q) = \frac{(q^{s-2\ell+1}-\varepsilon)(q^{s-2\ell}+\varepsilon)}{q^{2k_0-2}(q^{2\ell+2}-1)}.
\]
Assembling all three factors, the terms $(q^{s-2\ell+1}-\varepsilon)$ and
$(q^{s+1}-\varepsilon)$ cancel. Using $s+k_0 = n$,
$q^{2\ell+2}-1 = (q^{\ell+1}-1)(q^{\ell+1}+1)$, and the identity
$q^{s-2\ell}+\varepsilon = q^{b}-(-1)^{b}$
(which follows from $s-b = 2\ell$ and $\varepsilon = (-1)^{s+1} = (-1)^{b+1}$),
we obtain~\eqref{eq:alpha}.

Now we bound $\alpha^H$.
Set $N = q^{a+b-1}(q^{\ell+1}+1)$ and $D = (q^{a}-(-1)^{a})(q^{b}-(-1)^{b})$,
so that $\alpha^H = N/D$. We analyze $N - D$ by parity of the pair $(a,b)$.

\emph{Case $(a, b)$ both even.}
Here $D = q^{a+b}-q^{a}-q^{b}+1$, so
\[
  N - D = (q^{a+b+\ell}-q^{a+b}) + q^{a+b-1} + q^{a} + q^{b} - 1.
\]
The first term is non-negative ($q^{a+b+\ell} \geq q^{a+b}$), and since
$a, b \geq 2$ we have $q^{a+b-1} \geq q^3 \geq 8 > 1$, so $N - D > 0$.

\emph{Case $(a, b)$ of opposite parity.}
By symmetry assume $a$ odd, $b$ even. Then $D = q^{a+b}-q^{a}+q^{b}-1$ and
\[
  N - D = q^{a+b}(q^{\ell}-1) + q^{a+b-1} + q^{a} - q^{b} + 1.
\]
For $\ell\geq 1$, the first term satisfies $q^{a+b}(q^{\ell}-1) \geq q^{a+b}(q-1) \geq q^{a+b} > q^{b}$,
yielding $N-D > 0$.
For $\ell = 0$, $N - D = q^{b}(q^{a-1}-1) + q^{a} + 1$,
which is positive for $a \geq 2$, and equals $q+1 > 0$ when $a = 1$.

\emph{Case $(a, b)$ both odd.}
Here $D = q^{a+b}+q^{a}+q^{b}+1$ and
\[
  N - D = q^{a+b-1}(q^{\ell+1}+1) - q^{a+b} - q^{a} - q^{b} - 1.
\]
For $\ell \geq 1$, since $q^{\ell+1}+1 \geq q^{2}+1$, we have
\[
  N - D \geq q^{a+b-1}(q^{2}-q+1) - q^{a} - q^{b} - 1.
\]
Assuming without loss of generality that $a \leq b$, since $q^{2}-q+1 \geq 3$ and $a \geq 1$,
$3q^{a+b-1} \geq 3q^{b} > 2q^{b}+1 \geq q^{a}+q^{b}+1$, hence $N-D > 0$.

For $\ell = 0$, $N - D = q^{a+b-1} - q^{a} - q^{b} - 1$.
If $a \geq 3$ (so $b \geq 3$), $q^{a+b-1} = q^{a-1}q^{b} \geq q^{2}q^{b} \geq 4q^{b} > 2q^{b}+1 \geq q^{a}+q^{b}+1$.
If $a = 1$, $N - D = q^{b} - q - q^{b} - 1 = -(q+1) < 0$, the exceptional case.
Here
\[
  \alpha^H = \frac{q^{b}(q+1)}{(q+1)(q^{b}+1)} = \frac{q^{b}}{q^{b}+1},
\]
and $q^{b}/(q^{b}+1) \geq q/(q+1)$ for $b \geq 1$ (with equality at $b=1$).
The same bound holds by $a$-$b$ symmetry if $b = 1$.
Combining the cases, $\alpha^H > 1$ unless $\ell=0$, $(a,b)$ are both odd, and
$\min(a,b) = 1$, in which case $\alpha^H = q^{\max(a,b)}/(q^{\max(a,b)}+1) \in [q/(q+1), 1)$.
\end{proof}


\begin{corollary}\label{cor:Amonotone-strict}
For all $q \geq 2$ and $\ell+1\leq k\leq n-\ell-1$:
\[
  A^H_{n,k,\ell,q} > (q^{\ell+1}-1)\cdot A^H_{n,k,\ell+1,q}
\]
holds for all valid parameters except when $q = 2$, $\ell = 0$, $n$ is even, and $k\in\{1, n-1\}$.
In each such exceptional case,
$A^H_{n,k,0,2} = (2^{n-1}/(2^{n-1}+1))\cdot A^H_{n,k,1,2}$.
\end{corollary}

\begin{proof}
When $\alpha^H > 1$, the inequality immediately follows by~\eqref{eq:decomp}.
In the exceptional case of Theorem~\ref{thm:main},
$\alpha^H(q-1) = (q^{\max(a,b)}/(q^{\max(a,b)}+1))(q-1) > 1$
if and only if $(q-1)q^{\max(a,b)} > q^{\max(a,b)}+1$, which holds if and only if
$q\geq 3$. Hence $\alpha^H(q-1) \leq 1$ only when $q=2$, and in this case the
exceptional case of Theorem~\ref{thm:main} forces $\ell = 0$, both $a$ and $b$
odd, and $\min(a,b) = 1$; consequently $n = a + b$ is even, with $k = 1$ (when
$a = 1$) or $k = n - 1$ (when $b = 1$). For $q = 2$ the stated equality holds
by~\eqref{eq:alpha}.
\end{proof}

\begin{corollary}\label{cor:Amonotone}
For all $q \geq 2$ and $0\leq\ell<\min(k, n-k)$:
\[
  A^H_{n,k,\ell,q} > A^H_{n,k,\ell+1,q}
\]
holds for all valid parameters except $(q, \ell, k) = (2, 0, 1)$ and $(2, 0, n-1)$ with $n$ even.
\end{corollary}

\begin{proof}
Outside the exceptional family, $\alpha^H > 1$ and $q^{\ell+1}-1 \geq 1$, hence the product
exceeds~$1$ and the inequality holds.
In the exceptional family, $\alpha^H(q-1) = q^{n-1}/(q^{n-1}+1) \cdot 1 < 1$,
so $A^H_{n,k,0,2} < A^H_{n,k,1,2}$.
\end{proof}

\begin{corollary}[Asymptotic behavior of $A^H_\ell / A^H_{\ell+1}$]\label{cor:alpha-herm-asym}
Let $\ell \geq 0$ and $q$ be fixed, and set $a = k - \ell$, $b = n - k - \ell$
(so $a, b \geq 0$ and $a + b + 2\ell = n$). Recall $a$ measures the
``room'' above $\ell$ in the code dimension ($k = a + \ell$) and $b$ measures
the corresponding room in the codimension ($n - k = b + \ell$).
\begin{enumerate}
\item (\emph{Boundary regime: $k$ fixed, $n \to \infty$.})
Fix $a \geq 1$ and let $b \to \infty$. Then
\[
  \frac{A^H_{n,k,\ell,q}}{A^H_{n,k,\ell+1,q}}
  \;\longrightarrow\; \frac{q^{a-1}(q^{\ell+1}+1)(q^{\ell+1}-1)}{q^a - (-1)^a}.
\]
\item (\emph{Joint regime: $k \to \infty$ and $n - k \to \infty$.})
Let $a, b \to \infty$ simultaneously. Then
\[
  \frac{A^H_{n,k,\ell,q}}{A^H_{n,k,\ell+1,q}}
  \;\longrightarrow\; \frac{q^{2(\ell+1)} - 1}{q}.
\]
\item The asymptotic ratio in~(2) is minimized at $\ell = 0$, equal to
$(q^2 - 1)/q = q - q^{-1}$. This evaluates to $3/2$ at $q = 2$ and grows
like $q$ for large $q$, so $A^H_{n,k,\ell,q} > A^H_{n,k,\ell+1,q}$
asymptotically for every $\ell \geq 0$ and every $q \geq 2$.
\end{enumerate}
\end{corollary}

\begin{proof}
By Theorem~\ref{thm:main}, $A^H_\ell/A^H_{\ell+1} = \alpha^H \cdot (q^{\ell+1}-1)$,
and using $a + b + 2\ell = n$ the formula~\eqref{eq:alpha} can be rewritten as
$\alpha^H = q^{a+b-1}(q^{\ell+1}+1)/[(q^a-(-1)^a)(q^b-(-1)^b)]$.

(1) For fixed $a$ and $b \to \infty$, $(q^b - (-1)^b)/q^b \to 1$, so
\[
  \alpha^H
  = \frac{q^{a-1}(q^{\ell+1}+1)}{q^a - (-1)^a} \cdot \frac{q^b}{q^b - (-1)^b}
  \;\longrightarrow\; \frac{q^{a-1}(q^{\ell+1}+1)}{q^a - (-1)^a},
\]
and multiplying by $(q^{\ell+1}-1)$ gives the displayed formula.

(2) For $a, b \to \infty$ jointly, both $(q^a - (-1)^a)/q^a$ and
$(q^b - (-1)^b)/q^b$ tend to $1$, hence
$\alpha^H \to (q^{\ell+1}+1)/q$ and
$A^H_\ell/A^H_{\ell+1} \to (q^{\ell+1}+1)(q^{\ell+1}-1)/q = (q^{2(\ell+1)}-1)/q$.

(3) The map $\ell \mapsto (q^{2(\ell+1)}-1)/q$ is strictly increasing on
$\ell \geq 0$, so the minimum is attained at $\ell = 0$ with value
$(q^2 - 1)/q$. At $q = 2$ this equals $3/2 > 1$, and it increases without
bound in $q$, so the ratio strictly exceeds $1$ for all $q \geq 2$ and all
$\ell \geq 0$.
\end{proof}

\subsection{Numerical illustration of the boundary exception}\label{sec:Hnum}

Table~\ref{tab:hermitian-A} lists $A^H_{n,k,\ell,q}$ from Theorem~\ref{thm:LSL}
for small parameters at $q = 2$ (codes over $\F_4$) and $q = 3$ (codes over
$\F_9$). Boldface marks entries where $A^H_\ell$ strictly exceeds $A^H_{\ell-1}$,
confirming Corollary~\ref{cor:Amonotone-strict} on the boundary exception
family $\{(n, k, 0, 2) : n \text{ even},\ k \in \{1, n-1\}\}$.

\begin{table}[h]
\centering
\caption{$A^H_{n,k,\ell,q}$ (closed form, Theorem~\ref{thm:LSL}).
Boldface marks $A^H_{\ell+1} > A^H_\ell$ (boundary exception of
Corollary~\ref{cor:Amonotone-strict}).}
\label{tab:hermitian-A}
\begin{tabular}{ccc|rrrr}
\toprule
$n$ & $k$ & $q$ & $A^H_0$ & $A^H_1$ & $A^H_2$ & $A^H_3$ \\
\midrule
4 & 1 & 2 & 40 & \textbf{45} &  & \\
5 & 1 & 2 & 176 & 165 &  & \\
6 & 1 & 2 & 672 & \textbf{693} &  & \\
7 & 1 & 2 & 2752 & 2709 &  & \\
4 & 2 & 2 & 240 & 90 & 27 & \\
5 & 2 & 2 & 3520 & 1980 & 297 & \\
6 & 2 & 2 & 59136 & 27720 & 6237 & \\
6 & 3 & 2 & 197120 & 166320 & 12474 & 891 \\
7 & 2 & 2 & 924672 & 476784 & 89397 & \\
7 & 3 & 2 & 13561856 & 9535680 & 1072764 & 38313 \\
8 & 2 & 2 & 14970880 & 7368480 & 1519749 & \\
\midrule
4 & 1 & 3 & 540 & 280 &  & \\
5 & 1 & 3 & 4941 & 2440 &  & \\
6 & 1 & 3 & 44226 & 22204 &  & \\
4 & 2 & 3 & 5670 & 1680 & 112 & \\
5 & 2 & 3 & 444690 & 153720 & 6832 & \\
6 & 2 & 3 & 36420111 & 11990160 & 621712 & \\
6 & 3 & 3 & 312172380 & 125896680 & 3730272 & 27328 \\
\bottomrule
\end{tabular}
\end{table}

At $q = 2$, the bold entries at $[4,1]_4$ and $[6,1]_4$ confirm the
$A$-monotonicity boundary failure: precisely the family
$(q, \ell, k) = (2, 0, 1)$ or $(2, 0, n-1)$ with $n$ even. At $q = 3$ the
analogous parameters $[4,1]_9, [5,1]_9, [6,1]_9$ all satisfy
$A^H_0 > A^H_1$, consistent with Corollary~\ref{cor:Amonotone-strict}
restricting the failure to $q = 2$.

%% ============================================================
\section{Ratio decomposition for symplectic hull}\label{sec:symp}
%% ============================================================

We now establish the analogous theory for the symplectic inner product.
The structure parallels Section~\ref{sec:main}, but with two key differences:
the constraint $k - \ell \in 2\mathbb{Z}$ forces hull dimensions to step by $2$, and
the closed-form ratio $\alpha^S$ behaves qualitatively differently from $\alpha^H$.

\subsection{Main theorem}\label{sec:Smain}

Unlike the Hermitian case (Theorem~\ref{thm:LSL}), the symplectic mass formula
of Theorem~\ref{thm:LSL-symp} has no parity case-split: the formula is uniform
across all valid parameters. This makes the derivation of $\alpha^S$
algebraically more direct, and we can state the closed form and the
exception classification as a single result.

\begin{theorem}\label{thm:main-symp}
Let $n, k, \ell$ be non-negative integers with $\ell + 2 \leq k \leq 2n - \ell - 2$
and $k - \ell$ even. Set $a = k - \ell$ and $b = 2n - k - \ell$ (both even). Then
\[
  A^S_{2n,k,\ell,q} = \alpha^S_{2n,k,\ell,q} \cdot (q^{\ell+1}-1)(q^{\ell+2}-1) \cdot A^S_{2n,k,\ell+2,q},
\]
where
\begin{equation}\label{eq:alpha-symp}
  \alpha^S_{2n,k,\ell,q} = \frac{q^{a+b-2}}{(q^{a}-1)(q^{b}-1)}.
\end{equation}
\end{theorem}

\begin{proof}
Set $k_0 = (k-\ell)/2$. By Theorem~\ref{thm:LSL-symp}, we may write
$A^S_{2n, k, \ell, q} = q^{2k_0(n-k_0-\ell)} \cdot P_\ell \cdot \binom{n}{k_0}_{q^2}$
where $P_\ell = \prod_{m=1}^{\ell}(q^{2(n-k_0-\ell+m)}-1) / \prod_{m=1}^{\ell}(q^m-1)$.
Going from $\ell$ to $\ell+2$ sends $k_0 \to k_0 - 1$, so we compute the ratio
$A^S_\ell / A^S_{\ell+2}$ as a product of four factors arising from the four
ingredients of the formula:
\begin{itemize}
\item The prefactor ratio gives $q^{2(n-\ell-1)}$.
\item The numerator product ratio gives $1 / [(q^{2(n-k_0-\ell)}-1)(q^{2(n-k_0+1)}-1)]$.
\item The denominator product ratio gives $(q^{\ell+1}-1)(q^{\ell+2}-1)$.
\item The Gaussian binomial ratio gives $(q^{2(n-k_0+1)}-1) / (q^{2k_0}-1)$.
\end{itemize}
Multiplying these together, the factors $(q^{2(n-k_0+1)}-1)$ cancel. Using
$2(n-k_0-\ell) = b$ and $2k_0 = a$, the resulting ratio simplifies to
\[
  \frac{A^S_\ell}{A^S_{\ell+2}} = \frac{q^{a+b-2} (q^{\ell+1}-1)(q^{\ell+2}-1)}{(q^a-1)(q^b-1)},
\]
which gives~\eqref{eq:alpha-symp}.
\end{proof}

\medskip
The asymptotic ratio
$A^S_\ell / A^S_{\ell+2} \to (q^{\ell+1}-1)(q^{\ell+2}-1)/q^2$
(Corollary~\ref{cor:alpha-symp-asym}) falls below~$1$ at $\ell = 0$ when
$q = 2$ (where it equals $3/4$). Consequently, the natural BBO-style
monotonicity inequality $A^S_{\ell} > (q^{\ell+1}-1)(q^{\ell+2}-1) A^S_{\ell+2}$
fails for a substantial range of parameters at $q = 2$.

\begin{theorem}\label{thm:exception-symp}
For valid parameters $(2n, k, \ell, q)$:
\[
  A^S_{2n,k,\ell,q} > (q^{\ell+1}-1)(q^{\ell+2}-1) \cdot A^S_{2n,k,\ell+2,q}
\]
holds for \emph{all} $q \geq 3$, and for $q = 2$ if and only if $\ell \geq 1$
or $\min(k, 2n-k) \leq 2$.
The inequality fails precisely in the family
\[
  \mathcal{E}_S := \{(2n, k, 0, 2) : 4 \leq k \leq 2n-4,\ k \text{ even}\}.
\]
In each case in $\mathcal{E}_S$, we have $\alpha^S(q-1)(q^2-1) = 3 \cdot 2^{a+b-2}/[(2^a-1)(2^b-1)] < 1$
with $a = k$, $b = 2n - k$.
\end{theorem}

\begin{proof}
Set $N = q^{a+b-2}(q^{\ell+1}-1)(q^{\ell+2}-1)$ and $D = (q^a-1)(q^b-1)$.
We show $N > D$ except on $\mathcal{E}_S$.

\emph{Case $\ell \geq 1$.} Since $q^{\ell+1} \geq q^2$ and $q^{\ell+2} \geq q^3$,
\[
  (q^{\ell+1}-1)(q^{\ell+2}-1) \geq (q^2-1)(q^3-1) = q^5 - q^3 - q^2 + 1.
\]
Hence $N \geq q^{a+b-2}(q^5 - q^3 - q^2 + 1)$, and
\[
  N - D \geq q^{a+b}(q^3 - q - 2)/q + (q^{a+b-1} - q^a - q^b + 1).
\]
For $q \geq 2$, $q^3 - q - 2 \geq 4$, hence the first term is positive,
and the second term is positive for $a + b \geq 4$ (always satisfied).
Thus $N > D$.

\emph{Case $\ell = 0$ and $q \geq 3$.}
We need $q^{a+b-2}(q-1)(q^2-1) > (q^a-1)(q^b-1)$.
Using $(q^a-1)(q^b-1) < q^{a+b}$,
it suffices to show $(q-1)(q^2-1) > q^2$.
At $q = 3$, $(q-1)(q^2-1) = 2 \cdot 8 = 16 > 9 = q^2$,
and $(q-1)(q^2-1) - q^2 = q^3 - 2q^2 - q + 1$ is increasing for $q \geq 3$.
Hence the inequality holds.

\emph{Case $\ell = 0$ and $q = 2$.} The condition $N > D$ becomes
$3 \cdot 2^{a+b-2} > (2^a-1)(2^b-1)$, i.e.,
$2^{a+b-2} \cdot 3 > 2^{a+b} - 2^a - 2^b + 1$, equivalently
$2^a + 2^b > 2^{a+b-2} + 1$.
Without loss of generality assume $a \leq b$.
\begin{itemize}
\item If $a = 2$: $2^b + 4 > 2^b + 1$ holds trivially. (Monotone.)
\item If $a \geq 4$ (and $b \geq 4$): $2^a + 2^b \leq 2 \cdot 2^b \leq 2^{a+b-2}$ since $2^{a-2} \geq 4$, hence
$2^{a+b-2} \geq 4 \cdot 2^b \geq 2(2^a + 2^b)$, contradicting the required inequality.
(Non-monotone, $\mathcal{E}_S$ case.)
\end{itemize}
This completes the proof.
\end{proof}

\begin{corollary}[Asymptotic behavior of $A^S_\ell / A^S_{\ell+2}$]\label{cor:alpha-symp-asym}
Let $\ell \geq 0$ and $q$ be fixed, and set $a = k - \ell$, $b = 2n - k - \ell$
(so $a, b \geq 0$ and $a + b + 2\ell = 2n$). Recall $a$ measures the
``room'' above $\ell$ in the code dimension ($k = a + \ell$) and $b$ the
corresponding room in the codimension ($2n - k = b + \ell$).
\begin{enumerate}
\item (\emph{Boundary regime: $k$ fixed, $n \to \infty$.})
Fix $a \geq 2$ and let $b \to \infty$. Then
\[
  \frac{A^S_{2n,k,\ell,q}}{A^S_{2n,k,\ell+2,q}}
  \;\longrightarrow\; \frac{q^{a-2}(q^{\ell+1}-1)(q^{\ell+2}-1)}{q^a - 1}.
\]
\item (\emph{Joint regime: $k \to \infty$ and $2n - k \to \infty$.})
Let $a, b \to \infty$ simultaneously. Then
\[
  \frac{A^S_{2n,k,\ell,q}}{A^S_{2n,k,\ell+2,q}}
  \;\longrightarrow\; \frac{(q^{\ell+1}-1)(q^{\ell+2}-1)}{q^2}.
\]
\item The asymptotic ratio in~(2) is minimized at $\ell = 0$, equal to
$(q-1)(q^2-1)/q^2$. At $q = 2$ this equals $3/4 < 1$: the joint limit
$a, b \to \infty$ at $\ell = 0$, $q = 2$ is taken within the exception
family $\mathcal{E}_S$ of Theorem~\ref{thm:exception-symp}, and indeed
$\mathcal{E}_S$ contains arbitrarily large parameters, so the asymptotic
$3/4 < 1$ confirms that the BBO inequality fails throughout $\mathcal{E}_S$
as $2n$ grows. For $q \geq 3$, the asymptotic ratio is
$(q-1)^2(q+1)/q^2 \geq 16/9 > 1$, lying outside the exception family.
For each fixed $\ell \geq 1$ the asymptotic ratio is strictly larger (e.g.,
$A^S_1/A^S_3 \to 21/4$ at $q = 2$) and monotonicity holds asymptotically
in all cases.
\end{enumerate}
\end{corollary}

\begin{proof}
By Theorem~\ref{thm:main-symp},
$A^S_\ell/A^S_{\ell+2} = \alpha^S \cdot (q^{\ell+1}-1)(q^{\ell+2}-1)$ with
$\alpha^S = q^{a+b-2}/[(q^a-1)(q^b-1)]$.

(1) For fixed $a$ and $b \to \infty$, $(q^b - 1)/q^b \to 1$, so
\[
  \alpha^S = \frac{q^{a-2}}{q^a - 1} \cdot \frac{q^b}{q^b - 1}
  \;\longrightarrow\; \frac{q^{a-2}}{q^a - 1},
\]
and multiplying by $(q^{\ell+1}-1)(q^{\ell+2}-1)$ gives the displayed formula.

(2) For $a, b \to \infty$ jointly, both $(q^a - 1)/q^a$ and $(q^b - 1)/q^b$
tend to $1$, hence $\alpha^S \to q^{-2}$ and $A^S_\ell/A^S_{\ell+2} \to
(q^{\ell+1}-1)(q^{\ell+2}-1)/q^2$.

(3) The function $\ell \mapsto (q^{\ell+1}-1)(q^{\ell+2}-1)$ is strictly
increasing on $\ell \geq 0$, so the minimum over admissible $\ell$ is attained
at $\ell = 0$ with value $(q-1)(q^2-1)$. At $q = 2$ this divides by $q^2 = 4$
to give $3/4$. For $q \geq 3$, write $(q-1)(q^2-1) = (q-1)^2(q+1)$; the
function $(q-1)^2(q+1)/q^2$ is increasing in $q$ for $q \geq 3$, and at
$q = 3$ equals $4 \cdot 4/9 = 16/9$, so the quantity is $\geq 16/9$
throughout. At $\ell \geq 1$, the asymptotic ratio
$(q^{\ell+1}-1)(q^{\ell+2}-1)/q^2$ exceeds $(q^2-1)(q^3-1)/q^2 = 21/4 > 1$
at $q = 2, \ell = 1$, with strict monotone growth in $\ell$, hence exceeds
$1$ for all $\ell \geq 1$ and all $q \geq 2$.
\end{proof}

\subsection{Numerical illustration of the exception family
  \texorpdfstring{$\mathcal{E}_S$}{E\_S}}\label{sec:Snum}

Table~\ref{tab:symplectic-A} lists $A^S_{2n,k,\ell,q}$ from
Theorem~\ref{thm:LSL-symp} for small parameters. Boldface marks entries where
$A^S_\ell$ strictly exceeds $A^S_{\ell-2}$, confirming
Theorem~\ref{thm:exception-symp} on the family
$\mathcal{E}_S = \{(2n, k, 0, 2) : 4 \leq k \leq 2n - 4,\ k \text{ even}\}$
of $A$-monotonicity failures.

\begin{table}[h]
\centering
\caption{$A^S_{2n,k,\ell,q}$ (closed form, Theorem~\ref{thm:LSL-symp}).
Boldface marks $A^S_2 > A^S_0$.}
\label{tab:symplectic-A}
\begin{tabular}{ccc|rrrr}
\toprule
$2n$ & $k$ & $q$ & $A^S_0$ & $A^S_2$ & $A^S_4$ & $A^S_6$ \\
\midrule
4 & 2 & 2 & 20 & 15 & & \\
6 & 2 & 2 & 336 & 315 & & \\
8 & 2 & 2 & 5440 & 5355 & & \\
8 & 4 & 2 & 91392 & \textbf{107100} & 2295 & \\
10 & 2 & 2 & 87296 & 86955 & & \\
10 & 4 & 2 & 23744512 & \textbf{29216880} & 782595 & \\
12 & 4 & 2 & 6100942848 & \textbf{7596388800} & 213648435 & \\
12 & 6 & 2 & 98777169920 & \textbf{127619331840} & 4272968700 & 4922775 \\
\midrule
4 & 2 & 3 & 90 & 40 & & \\
6 & 2 & 3 & 7371 & 3640 & & \\
8 & 2 & 3 & 597780 & 298480 & & \\
8 & 4 & 3 & 48958182 & 26863200 & 91840 & \\
\bottomrule
\end{tabular}
\end{table}

At $q = 2$, all entries with $4 \leq k \leq 2n - 4$ exhibit
$A^S_0 < A^S_2$, consistent with Theorem~\ref{thm:exception-symp}.
At $q = 3$, $A^S$ is strictly monotone in $\ell$ in every computed case,
consistent with Theorem~\ref{thm:exception-symp} restricting failures to $q = 2$.

%% ============================================================
\section{Comparison and applications}\label{sec:comp-app}
%% ============================================================

Table~\ref{tab:comparison} juxtaposes the Euclidean (BBO~\cite{BBOz2026}),
Hermitian (Section~\ref{sec:main}), and symplectic (Section~\ref{sec:symp}) ratio
decompositions side by side.

\begin{table}[h]
\centering
\footnotesize
\caption{Comparison of the ratio decompositions across the three classical
inner products. Here $a = k - \ell$, $b = n - k - \ell$ (Euclidean/Hermitian)
or $b = 2n - k - \ell$ (symplectic); $\Delta$ is the step in $\ell$.}
\label{tab:comparison}
\setlength{\tabcolsep}{2pt}
\begin{tabular}{l|c|c|c}
\toprule
& Euclidean~\cite{BBOz2026} & Hermitian & Symplectic \\
\midrule
Step $\Delta$ in $\ell$ & $1$ & $1$ & $2$ \\[2pt]
$\alpha^*$ closed form &
   case-split &
   $\dfrac{q^{n-2\ell-1}(q^{\ell+1}+1)}{(q^a-(-1)^a)(q^b-(-1)^b)}$ &
   $\dfrac{q^{a+b-2}}{(q^a-1)(q^b-1)}$ \\[6pt]
$\alpha^*$ lower bound & $1/2$ & $q/(q+1) \geq 2/3$ & none $> 1$ \\
$\alpha^*$ asymptotic & $(q+1)/q$ & $(q+1)/q$ & $1/q^2$ \\[2pt]
Asymp.\ $A_0/A_\Delta$ & $(q^2{-}1)/q$ & $(q^2{-}1)/q$ & $(q{-}1)(q^2{-}1)/q^2$ \\
\;at $q = 2$ & $3/2$ & $3/2$ & $\mathbf{3/4}$ \\
\;at $q \geq 3$ & $\geq 8/3$ & $\geq 8/3$ & $\geq 16/9$ \\
\midrule
Exceptions & $\eta$-char + $\bmod 4$ & $\ell{=}0,\,n$ even, $k{\in}\{1,n{-}1\}$ & $q{=}2,\,\ell{=}0,\,4{\leq}k{\leq}2n{-}4$ \\
\bottomrule
\end{tabular}
\end{table}

The contrast in Table~\ref{tab:comparison} tracks the Witt classification
of the underlying form. The orthogonal group over odd $\F_q$ splits
non-degenerate subspaces into two isometry classes (the discriminant lying in
$\F_q^*/(\F_q^*)^2$), producing the $\eta$-character case-split in the BBO
formula and the worst-case bound $\alpha \geq 1/2$. The unitary and symplectic
groups, in contrast, each act with a single orbit on LCD codes, so the closed
forms $\alpha^H$ and $\alpha^S$ contain no quadratic-character factor. The
symplectic case further degenerates: every vector is self-orthogonal, the
parity constraint $k - \ell \in 2\mathbb{Z}$ is forced, and the loss of a
hyperbolic pair per step of~$2$ in $\ell$ yields an asymptotic ratio of
$1/q^2$ rather than $(q+1)/q$. Consequently, the two exception sets to the
BBO-style inequality are qualitatively different: a $0$-dimensional Hermitian
boundary versus a $1$-parameter symplectic family of generic failures at
$q = 2$.

The closed-form counts $A^H$ and $A^S$ have direct relevance to quantum
coding theory, sketched below.

\subsection{Symplectic entanglement-assisted stabilizer codes}

Relaxing the symplectic self-orthogonality required by KKKS-style
stabilizer constructions~\cite{KKKS2006} to allow shared entanglement
yields the entanglement-assisted stabilizer formalism of Brun, Devetak,
and Hsieh~\cite{BDH2006}. The minimum number of pre-shared maximally
entangled pairs needed by such a construction was determined by Wilde
and Brun~\cite[Theorem~1]{WildeBrun2008}.

\begin{theorem}[{\cite[Theorem~1]{WildeBrun2008}}]\label{thm:WB}
Let $H = [H_Z \mid H_X]$ be an $(n - k) \times 2n$ binary quantum check
matrix. The resulting entanglement-assisted stabilizer code has parameters
$[[n, k + c; c]]_2$, where the optimal number of ebits is
\[
   c \;=\; \tfrac{1}{2}\,\mathrm{rank}\bigl(H_X H_Z^T + H_Z H_X^T\bigr).
\]
\end{theorem}

The matrix $H_X H_Z^T + H_Z H_X^T$ is the Gram matrix of the rows of $H$
under the symplectic form $\langle\cdot,\cdot\rangle_S$ of
Section~\ref{sec:prelim-dual}; viewing the row span of $H$ as an
$\F_2$-linear $[2n, m]_2$ code $C = \langle H\rangle$ of dimension
$m = n - k$, the Gram rank equals $m - \dim(\mathrm{Hull}_S(C))$ and
hence
\[
   c \;=\; \tfrac{1}{2}\bigl(m - \dim(\mathrm{Hull}_S(C))\bigr).
\]
An analogous formula holds for $\F_q$-linear codes via the

\begin{thebibliography}{99}
\bibitem[BBOz2026]{BBOz2026} S.~Bouyuklieva, I.~Bouyukliev, F.~\"Ozbudak, Sequence of numbers of linear codes with increasing hull dimensions, \textit{Des.\ Codes Cryptogr.} \textbf{94} (2026). \url{https://doi.org/10.1007/s10623-025-01776-9}.
\bibitem[BDH2006]{BDH2006} T.~A.~Brun, I.~Devetak, M.-H.~Hsieh, Correcting quantum errors with entanglement, \textit{Science} \textbf{314} (2006), 436--439.
\bibitem[CarletGuilley2016]{CarletGuilley2016} C.~Carlet, S.~Guilley, Complementary dual codes for counter-measures to side-channel attacks, \textit{Adv.\ Math.\ Commun.} \textbf{10} (2016), 131--150.
\bibitem[Dcc-EAqecc]{Dcc-EAqecc} K.~Guenda, S.~Jitman, T.~A.~Gulliver, Constructions of good entanglement-assisted quantum error correcting codes, \textit{Des.\ Codes Cryptogr.} \textbf{86}(1) (2018), 121--136.
\bibitem[KKKS2006]{KKKS2006} A.~Ketkar, A.~Klappenecker, S.~Kumar, P.~K.~Sarvepalli, Nonbinary stabilizer codes over finite fields, \textit{IEEE Trans.\ Inform.\ Theory} \textbf{52}(11) (2006), 4892--4914.
\bibitem[LSL2024]{LSL2024} S.~Li, M.~Shi, Y.~Li, S.~Ling, A further study on the mass formula for linear codes with prescribed hull dimension, arXiv:2410.13578, 2024.
\bibitem[Massey1992]{Massey1992} J.~L.~Massey, Linear codes with complementary duals, \textit{Discrete Math.} \textbf{106/107} (1992), 337--342.
\bibitem[WildeBrun2008]{WildeBrun2008} M.~M.~Wilde, T.~A.~Brun, Optimal entanglement formulas for entanglement-assisted quantum coding, \textit{Phys.\ Rev.\ A} \textbf{77} (2008), 064302.
\end{thebibliography}
\end{document}
